\documentclass[11pt,a4paper]{article}
\usepackage[utf8]{inputenc}
\usepackage[T1]{fontenc}
\usepackage{lmodern}
\usepackage[italian]{babel}
\usepackage{amsmath,amssymb}
\usepackage[margin=2.4cm]{geometry}
\usepackage{booktabs}
\usepackage{hyperref}

\newcommand{\R}{\mathbb{R}}
\newcommand{\E}{\mathbb{E}}
\newcommand{\norm}[1]{\left\lVert #1 \right\rVert}
\newcommand{\fonte}[1]{\par\noindent\textit{Fonte: #1}\par\smallskip}
\newcommand{\esterna}[1]{\par\noindent\textit{Fonte esterna (non nei materiali del corso): #1}\par\smallskip}

\title{Double descent nei modelli lineari\\\large Teoria per l'Assignment 4 (SMML 2025/26, prof.\ Cesa-Bianchi)}
\author{Filippo Ferroni}
\date{Versione 4, 25 ottobre 2026}

\begin{document}
\maketitle

\paragraph{Come usare questo documento.}
Leggilo in circa 40 minuti. Le sezioni 2--6 sono la teoria minima da saper spiegare all'orale.
La sezione 7 collega la teoria all'estensione con la gradient descent, la sezione 8 descrive cosa ci aspettiamo
di vedere nei grafici, la sezione 9 il disegno degli esperimenti. In fondo: domande d'orale con risposta e soluzioni
delle domande di verifica. Per ogni affermazione indico la fonte; quando un concetto \textbf{non} è nei materiali del corso
del corso (slide 1--15 di Cesa-Bianchi) lo segnalo esplicitamente come fonte esterna.

\tableofcontents

%--------------------------------------------------------------------
\section{Notazione e corrispondenze}

\begin{center}
\begin{tabular}{@{}llll@{}}
\toprule
Concetto & Corso (slide 9) & Assignment / paper & Codice \\
\midrule
numero di esempi di training & $m$ & $n$ & \texttt{m} \\
numero di feature (complessità) & $d$ & $d$ (assignment) / $N$ (paper) & \texttt{d} \\
design matrix & $X \in \R^{m\times d}$ & & \texttt{X} \\
etichette & $y \in \R^m$ & & \texttt{y} \\
parametro ridge & $\alpha$ & $\lambda$ (altri testi) & \texttt{alpha} \\
soglia di interpolazione & $d = m$ & $N = n$ & \\
\bottomrule
\end{tabular}
\end{center}

\textbf{Attenzione:} nel paper $d$ è la dimensione dell'\emph{input} $x$, mentre il numero di feature/parametri è $N$ e i campioni sono $n$ (p.~2 e p.~4).

Nel report scriveremo una volta sola: ``seguiamo la notazione del corso; il numero di esempi $n$ dell'assignment
corrisponde a $m$''. Le righe di $X$ sono gli esempi: $X^\top = [x_1,\dots,x_m]$ (slide 9, p.~5).
Un predittore lineare per la regressione è $h(x) = w^\top x$ con $w \in \R^d$ (slide 9, p.~1 e p.~4).
L'intercetta si gestisce aggiungendo una feature costante pari a $1$ (slide 9, p.~1) oppure centrando i dati.

%--------------------------------------------------------------------
\section{Setup statistico: rischio, errore di training, errore di test}
\fonte{slide ``4. statRisk'' (p.~1--4) e ``1. intro'' (p.~3--5).}

\begin{itemize}
  \item Ogni esempio $(x,y)$ è estratto in modo indipendente da una distribuzione $\mathcal{D}$ fissata ma sconosciuta
        (statRisk, p.~1). Training set e test set sono quindi campioni statistici.
  \item \textbf{Rischio statistico} di un predittore $h$ rispetto alla loss $\ell$ (statRisk, p.~1):
        $\ell_{\mathcal{D}}(h) = \E\bigl[\ell(h(X), Y)\bigr]$. Per noi $\ell(\hat y, y) = (\hat y - y)^2$ e $h(x) = w^\top x$.
  \item \textbf{Predittore Bayes-ottimo} $f^\ast$: il migliore possibile data $\mathcal{D}$. Per la square loss
        $f^\ast(x) = \E[Y \mid X = x]$ e il \textbf{rischio di Bayes} è la varianza condizionata attesa dell'etichetta
        (statRisk, p.~1--2). Nel nostro dataset sintetico, con $y = w^{\star\top}x + \varepsilon$ e rumore di varianza
        $\sigma^2$, il rischio di Bayes è proprio $\sigma^2$: \textbf{nessun modello può avere test error atteso sotto $\sigma^2$}.
        Lo useremo come linea di riferimento nei grafici.
  \item \textbf{Errore di training}: $\ell_S(h) = \frac{1}{m}\sum_{t=1}^{m} \ell(h(x_t), y_t)$; per noi
        $\ell_S(w) = \frac{1}{m}\norm{Xw - y}^2$. ERM sceglie il predittore che lo minimizza (intro, p.~4).
  \item \textbf{Errore di test}: la stessa media calcolata su un test set $S'$ indipendente; stima il rischio
        $\ell_{\mathcal{D}}(h)$, e per Chernoff--Hoeffding la stima è accurata quando il test set è grande (statRisk, p.~3--4).
        Il test set non va mai usato per addestrare né per scegliere iperparametri.
\end{itemize}

\paragraph{Overfitting e underfitting nel corso.}
Intro, p.~5: se l'algoritmo restituisce un predittore con training error alto è \emph{underfitting}; con training error
basso ma test error alto è \emph{overfitting}. Il corso osserva che ci si aspetta overfitting quando la classe di
predittori è troppo ricca rispetto a $m$ (``$\log_2|F| \gg m$'', intro p.~5). \textbf{Il double descent è esattamente un caso
in cui questa intuizione classica non basta}: per $d \gg m$ il training error è $0$ ma il test error torna a scendere.

%--------------------------------------------------------------------
\section{Least squares con $d \le m$: la formula delle slide}
\fonte{slide 9, p.~4--5.}

Il predittore di regressione lineare è l'ERM rispetto alla square loss:
\[
  w_S = \operatorname*{argmin}_{w \in \R^d} \norm{Xw - y}^2 .
\]
$F(w) = \norm{Xw - y}^2$ è convessa, quindi il minimo si trova dove il gradiente è nullo:
\[
  \nabla F(w) = 2X^\top(Xw - y) = 0
  \quad\Longleftrightarrow\quad
  X^\top X\, w = X^\top y \qquad \text{(equazioni normali)}.
\]
Se $X^\top X$ è invertibile:
\[
  \boxed{w_S = (X^\top X)^{-1} X^\top y}
\]
\textbf{Quando è invertibile?} La slide dice: se e solo se $x_1,\dots,x_m$ generano $\R^d$.
Servono almeno $d$ vettori per generare $\R^d$, quindi serve $m \ge d$.

\paragraph{Interpretazione geometrica.}
Il vettore delle predizioni sul training set è $Xw_S = X(X^\top X)^{-1}X^\top y = Py$, dove $P$ è la proiezione
ortogonale sullo spazio generato dalle colonne di $X$. Più feature aggiungiamo, più grande è quello spazio, più
$Py$ si avvicina a $y$: l'errore di training \textbf{non può aumentare} quando $d$ cresce.

\paragraph{Il caso $d = m$.}
Se $X$ è quadrata e invertibile, $w_S = X^{-1}y$ e $Xw_S = y$: errore di training esattamente $0$.
Il modello \textbf{interpola} già a $d = m$. Da qui il nome ``soglia di interpolazione''.

%--------------------------------------------------------------------
\section{Least squares con $d > m$: la soluzione a minima norma}
\esterna{pseudo-inversa di Moore--Penrose, testi standard di algebra lineare (es.\ Strang, \emph{Introduction to Linear
Algebra}, capitolo su pseudo-inversa e SVD); uso della soluzione a minima norma: Belkin et al.\ 2019, sez.~2 (p.~4: ``quando il minimizzatore non è unico \dots{} scegliamo quello con norma $\ell_2$ minima'') e p.~10 (``minimum norm linear regression in the feature space'').}

\subsection{Perché la formula delle slide non funziona più}
$X^\top X$ è una matrice $d \times d$, ma il suo rango è uguale al rango di $X$, che è al massimo $\min(m,d) = m$.
Se $d > m$ allora $\operatorname{rank}(X^\top X) \le m < d$: la matrice è \textbf{singolare} e $(X^\top X)^{-1}$ non esiste.
Equivalentemente: $m$ vettori non possono generare $\R^d$ se $m < d$.

\subsection{Infinite soluzioni che interpolano}
Con $d > m$ (e righe di $X$ linearmente indipendenti, il caso tipico con dati continui), il sistema $Xw = y$ ha
$m$ equazioni e $d$ incognite: ha \textbf{infinite} soluzioni esatte. Se $w_0$ è una soluzione, lo è anche $w_0 + z$
per ogni $z$ nel nucleo di $X$ ($Xz = 0$), che ha dimensione $d - m$. Tutte hanno errore di training $0$.

\subsection{Quale scegliere: la soluzione a minima norma}
\[
  \boxed{w^{+} = \operatorname*{argmin}\bigl\{\norm{w} : Xw = y\bigr\} = X^\top (X X^\top)^{-1} y}
\]
\textbf{Verifica che interpola:} $Xw^{+} = XX^\top(XX^\top)^{-1}y = y$. ($XX^\top$ è $m\times m$ e invertibile se
le righe di $X$ sono indipendenti.)

\textbf{Perché ha norma minima (dimostrazione breve, da saper fare).}
$w^{+} = X^\top a$ con $a = (XX^\top)^{-1}y$, quindi $w^{+}$ è combinazione lineare delle righe di $X$.
Ogni altra soluzione è $w = w^{+} + z$ con $Xz = 0$, cioè $z$ è ortogonale a tutte le righe di $X$, e quindi a $w^{+}$.
Per Pitagora:
\[
  \norm{w}^2 = \norm{w^{+}}^2 + \norm{z}^2 \ge \norm{w^{+}}^2 .
\]

\subsection{Una formula unica: la pseudo-inversa}
La pseudo-inversa $X^{+}$ si calcola dalla SVD $X = U\Sigma V^\top$: $X^{+} = V\Sigma^{+}U^\top$, dove $\Sigma^{+}$
inverte i valori singolari non nulli e lascia a zero gli altri. Vale:
\begin{itemize}
  \item se $d \le m$ e $X$ ha rango pieno: $X^{+}y = (X^\top X)^{-1}X^\top y = w_S$ (la formula delle slide);
  \item se $d > m$ e $X$ ha rango $m$: $X^{+}y = X^\top(XX^\top)^{-1}y = w^{+}$.
\end{itemize}
Quindi in numpy basta \texttt{w = np.linalg.pinv(X) @ y} per \textbf{tutti} i $d$.
È algebra lineare, non un ``fit'' di libreria, quindi rispetta il vincolo ``from scratch'' (FAQ dell'assignment, p.~16).

\paragraph{Note numeriche (importanti per il codice).}
\begin{itemize}
  \item Non usare \texttt{np.linalg.inv(X.T @ X)}: formare $X^\top X$ eleva al quadrato il numero di condizionamento,
        e vicino a $d=m$ i risultati diventano numericamente inaffidabili.
  \item \texttt{np.linalg.pinv} azzera solo i valori singolari sotto una soglia minuscola (\texttt{rcond}, circa $10^{-15}$
        relativo al più grande). Vicino a $d=m$ i valori singolari piccoli \emph{non} vengono azzerati: \texttt{pinv}
        non regolarizza, quindi il picco resta visibile. È il comportamento che vogliamo mostrare.
\end{itemize}

\subsection{Esempio numerico minimo}
$m=1$, $d=2$: un solo esempio $x = (1,1)$, $y = 2$. Le soluzioni di $w_1 + w_2 = 2$ sono infinite: $(2,0)$, $(0,2)$,
$(1,1)$, $(3,-1)$\dots{} La soluzione a minima norma è $w^{+} = x^\top(xx^\top)^{-1}y = (1,1)\cdot\frac{2}{2} = (1,1)$,
con norma $\sqrt{2}$, contro norma $2$ di $(2,0)$. Distribuisce il peso in modo uniforme invece di
``puntare tutto'' su una feature.

%--------------------------------------------------------------------
\section{Perché il test error esplode a $d \approx m$}
\fonte{slide 9, p.~5: ``quando $X^\top X$ è quasi singolare, $w_S$ è molto sensibile a perturbazioni del training set.
Questa instabilità aumenta l'errore di stima (varianza)''. È l'aggancio al corso.}

\subsection{Intuizione con un esempio $2\times 2$}
$m = d = 2$, righe $x_1 = (1,0)$, $x_2 = (1,\varepsilon)$ con $\varepsilon$ piccolo (righe quasi parallele, $X$ quasi
singolare). Etichette $y = (1,\, 1+\delta)$, dove $\delta$ è un piccolo rumore.
Risolvendo $Xw = y$: $w_1 = 1$ e $w_1 + \varepsilon w_2 = 1+\delta$, quindi
\[ w_2 = \delta/\varepsilon . \]
Con $\varepsilon = 0.01$ e un rumore $\delta = 0.1$, otteniamo $w_2 = 10$: un rumore piccolo produce un peso enorme.
Su un nuovo punto, ad esempio $x = (0,1)$, la predizione vale $10$ invece di circa $0$. È esattamente il picco.

\subsection{Versione generale con gli autovalori}
Siano $\lambda_1 \ge \dots \ge \lambda_d \ge 0$ gli autovalori di $X^\top X$ e $v_1,\dots,v_d$ gli autovettori
(ortonormali). La soluzione a minima norma si scrive
\[
  w^{+} = \sum_{i:\,\lambda_i > 0} \frac{v_i^\top X^\top y}{\lambda_i}\, v_i .
\]
Supponiamo label lineari con rumore: $y = Xw^\star + \varepsilon$. La parte dovuta al rumore è
$\sum_i \frac{v_i^\top X^\top \varepsilon}{\lambda_i} v_i$; ogni termine ha norma
$\frac{|u_i^\top \varepsilon|}{\sqrt{\lambda_i}}$, con $u_i = Xv_i/\sqrt{\lambda_i}$ vettore unitario.
\textbf{Il rumore viene diviso per $\sqrt{\lambda_i}$.}
\begin{itemize}
  \item $d \ll m$: $X$ è ``alta e stretta'', tutti gli autovalori sono grandi (circa $m$ per feature standardizzate):
        il rumore viene smorzato.
  \item $d \approx m$: $X$ è quasi quadrata e casuale, il più piccolo autovalore è vicinissimo a $0$:
        il rumore viene amplificato, $\norm{w}$ esplode, e con lui il test error.
  \item $d \gg m$: $X$ è ``larga'', gli $m$ autovalori non nulli tornano lontani da $0$ e il rumore viene ``spalmato'' su
        molte direzioni: $\norm{w^{+}}$ e il test error scendono. È la \textbf{seconda discesa}.
\end{itemize}
\esterna{il comportamento del più piccolo autovalore di una matrice casuale in funzione di $d/m$ è un risultato di random
matrix theory (legge di Marchenko--Pastur). Non lo dimostriamo: lo mostriamo empiricamente plottando $\norm{w}$
in funzione di $d$, che è più difendibile all'orale.}

\subsection{Una formula esatta per il regime classico (facoltativa)}
\esterna{Hastie, Montanari, Rosset, Tibshirani, ``Surprises in high-dimensional ridgeless least squares interpolation'',
Annals of Statistics 2022.}
Se $x \sim \mathcal{N}(0, I_d)$, $y = w^{\star\top}x + \varepsilon$ con rumore di varianza $\sigma^2$ e $d < m-1$, allora
\[
  \E\bigl[\norm{w_S - w^\star}^2\bigr] = \sigma^2\,\frac{d}{m - d - 1},
\]
che \textbf{diverge per $d \to m-1$}. Non serve dimostrarla; è utile per dire all'orale che il picco non è un artefatto
numerico ma un fenomeno statistico previsto dalla teoria.

%--------------------------------------------------------------------
\section{Ridge regression}
\fonte{slide 9, p.~5.}

\subsection{Formula}
\[
  w_{S,\alpha} = \operatorname*{argmin}_{w \in \R^d} \norm{Xw - y}^2 + \alpha \norm{w}^2 ,
  \qquad \alpha > 0 .
\]
Gradiente: $2X^\top(Xw - y) + 2\alpha w = 0 \;\Rightarrow\; (X^\top X + \alpha I)\,w = X^\top y$, quindi
\[
  \boxed{w_{S,\alpha} = (\alpha I + X^\top X)^{-1} X^\top y}
\]

\subsection{Perché elimina il picco}
Gli autovalori di $\alpha I + X^\top X$ sono $\alpha + \lambda_i \ge \alpha > 0$ (slide 9, p.~5): la matrice è sempre
invertibile, per ogni $d$, anche $d > m$. Nella base degli autovettori:
\[
  w_{S,\alpha} = \sum_{i} \frac{v_i^\top X^\top y}{\lambda_i + \alpha}\, v_i .
\]
Si divide per $\lambda_i + \alpha$ invece che per $\lambda_i$: anche quando $\lambda_i \approx 0$ il denominatore resta
almeno $\alpha$, quindi il rumore non può essere amplificato senza limite. Il picco si attenua o sparisce.

\subsection{Il prezzo: bias}
\begin{itemize}
  \item $\alpha \to 0$: si ritrova la soluzione a minima norma $w^{+}$ (anche per $d > m$; nelle direzioni con
        $\lambda_i = 0$ il numeratore $v_i^\top X^\top y$ è $0$). \emph{Fatto standard, non nelle slide.}
  \item $\alpha \to \infty$: $w_{S,\alpha} \to 0$, il predittore ignora i dati (slide 9, p.~5: $\alpha$ controlla il bias).
  \item In mezzo: $\alpha$ scambia varianza (instabilità) con bias (sotto-adattamento). Con $\alpha$ troppo grande il test
        error sale \emph{per tutti} i $d$.
\end{itemize}
\fonte{Belkin et al.\ 2019, p.~10: la regolarizzazione ``può sia impedire l'interpolazione sia cambiare la capacità effettiva della classe, attenuando o mascherando il picco''.}
\esterna{Nakkiran, Venkat, Kakade, Ma, ``Optimal regularization can mitigate double descent'', ICLR 2021: con $\alpha$
ottimale la curva del test error diventa monotona.}

\subsection{Calcolo in pratica}
\begin{itemize}
  \item Primale: \texttt{np.linalg.solve(X.T @ X + alpha * np.eye(d), X.T @ y)} (sistema $d \times d$).
  \item Duale, più economico per $d > m$: $w_{S,\alpha} = X^\top(XX^\top + \alpha I)^{-1}y$ (sistema $m\times m$).
        Le due forme coincidono perché $X^\top(XX^\top + \alpha I) = (X^\top X + \alpha I)X^\top$.
  \item \texttt{solve} è preferibile a calcolare l'inversa esplicita (più stabile e più veloce).
\end{itemize}

\subsection{Come scegliere $\alpha$ senza data leakage}
\begin{itemize}
  \item Per il grafico principale: alcune curve con $\alpha$ fissi (es.\ $10^{-3}, 10^{-1}, 10$) accanto alla curva
        di minima norma. Nessuna scelta sui dati, nessun leakage.
  \item Se avanza tempo: $\alpha$ scelto per ogni $d$ con un validation split (o $K$-fold cross-validation) ricavato
        \textbf{solo} dal training set (slide ``6. crossVal'', p.~2--3: tuning su validation set; p.~3--4: nested CV),
        come nel notebook \texttt{CrossValidation.ipynb} del prof.
        Il test set non si usa mai per scegliere.
\end{itemize}

%--------------------------------------------------------------------
\section{Collegamento con la gradient descent (estensione)}
\fonte{slide 10 (online gradient descent) per l'algoritmo; il risultato sulla minima norma è esterno (vedi sotto).}

Gradient descent su $F(w) = \norm{Xw - y}^2$ partendo da $w_1 = 0$:
\[
  w_{k+1} = w_k - \eta\, 2X^\top(Xw_k - y).
\]
\begin{itemize}
  \item Ogni aggiornamento è $X^\top(\dots)$, cioè una combinazione lineare delle righe di $X$. Partendo da $0$,
        $w_k$ resta sempre nello spazio delle righe.
  \item Con passo $\eta$ abbastanza piccolo ($\eta < 1/\lambda_1$), per $d > m$ la GD converge a una soluzione che
        interpola \emph{e} sta nello spazio delle righe: per la dimostrazione della Sezione 4.3, è proprio $w^{+}$.
  \item Fermarsi prima della convergenza (early stopping) ha un effetto simile alla ridge: le direzioni con $\lambda_i$
        piccolo sono le ultime a essere apprese, quindi il rumore lungo di esse non viene amplificato.
\end{itemize}
\fonte{bias implicito della GD verso la minima norma: Belkin et al.\ 2019, p.~11 (``SGD initialized at zero also converges to the minimum norm solution'').}
L'estensione si fa \textbf{solo} se report e repo sono completi.

%--------------------------------------------------------------------
\section{Cosa ci aspettiamo nei grafici}

\subsection{La curva a U classica e il double descent}
\begin{itemize}
  \item \textbf{Regime classico} ($d < m$): il corso scompone il rischio di $h_S$ in
        \emph{errore di stima} (varianza) $+$ \emph{errore di approssimazione} (bias) $+$ rischio di Bayes
        (statRisk, p.~4). Aggiungendo feature il bias scende ma la varianza sale; overfitting quando domina la varianza,
        underfitting quando domina il bias (statRisk, p.~6). Vicino a $d = m$ la varianza va a infinito: è la curva a U.
  \item \textbf{Soglia di interpolazione} ($d = m$): l'errore di training arriva a $0$, il test error ha il picco.
  \item \textbf{Regime sovra-parametrizzato} ($d > m$): errore di training sempre $0$, il test error scende di nuovo e
        può andare sotto il minimo del regime classico (paper, Fig.~1, p.~3). Spiegazione del paper (p.~3 e p.~6): oltre la soglia, classi più ricche contengono interpolanti di norma più piccola, cioè ``più semplici'' (rasoio di Occam); la norma è massima alla soglia e poi decresce (Fig.~2, p.~5). La teoria classica, che prevede solo la U,
        non spiega questa parte: è il punto del paper.
\end{itemize}

\subsection{Grafici previsti}
\begin{enumerate}
  \item Errore di training e di test in funzione di $d$ (minima norma), linea verticale a $d = m$, asse $y$ logaritmico
        (il picco è di ordini di grandezza).
  \item Sul sintetico, linea orizzontale al rischio di Bayes $\sigma^2$ (Sezione 2) come riferimento.
  \item $\norm{w}$ in funzione di $d$: deve avere il picco nello stesso punto (evidenza del meccanismo della Sezione 5).
  \item Test error di minima norma e ridge (più valori di $\alpha$) sullo stesso grafico.
  \item Gli stessi grafici per il dataset reale, e un confronto tra i due dataset.
\end{enumerate}

%--------------------------------------------------------------------
\section{Disegno degli esperimenti (anteprima)}

\paragraph{Rapporto con il paper.} Il paper mostra il fenomeno con random Fourier features e random ReLU features (sez.~2 e App.~C.1--C.2), reti neurali (sez.~2) e foreste/boosting (sez.~3), su MNIST e altri dataset di classificazione. Non contiene un esperimento di regressione lineare ``pura'', ma le random features sono addestrate con regressione lineare a minima norma nello spazio delle feature (p.~10): il nostro esperimento ne è la versione più semplice.

\subsection{Dataset sintetico: label lineari rumorose}
\esterna{modello a ``feature deboli'': Belkin, Hsu, Xu, ``Two models of double descent for weak features'', 2019.}
\begin{itemize}
  \item $x \sim \mathcal{N}(0, I_D)$ con $D$ grande (es.\ $D = 1000$), $y = w^{\star\top}x + \varepsilon$, rumore gaussiano.
  \item Il modello con complessità $d$ usa solo le prime $d$ feature. Per $d < D$ le feature non usate agiscono come
        rumore aggiuntivo (bias); per $d$ grande il modello ha accesso a tutto il segnale.
  \item $m$ fisso (es.\ $m = 100$), $d$ da $1$ a $D$, griglia più fitta intorno a $d = m$.
  \item Media (o mediana) su più ripetizioni con seed diversi, con banda di variabilità.
\end{itemize}

\subsection{Dataset reale con feature casuali}
\begin{itemize}
  \item Un dataset di regressione con poche feature (es.\ California Housing, 8 feature, presente anche nel repo del prof).
  \item Per far crescere $d$ oltre $m$: feature casuali $\phi(x) = \max(0, Wx)$ (ReLU) con $W$ gaussiana fissata,
        oppure random Fourier features come nel paper. Il modello resta \textbf{lineare nei parametri} $w$.
        Il corso chiama questa idea \emph{feature expansion}: aggiungere nuove feature costruite dalle originali per
        ridurre il bias dei predittori lineari (slide ``11. kernels'', p.~1).
  \item Sottocampioniamo $m$ esempi di training (es.\ $m = 200$) e teniamo un test set separato.
\end{itemize}

\subsection{Regole contro il data leakage}
\begin{itemize}
  \item Split train/test \emph{prima} di qualunque pre-processing.
  \item Media e deviazione standard per standardizzare calcolate \textbf{solo} sul training set e applicate al test set.
  \item Le matrici casuali $W$ non dipendono dai dati, quindi non creano leakage, ma vanno generate con seed fissato.
  \item Iperparametri (es.\ $\alpha$) scelti solo con dati di training.
\end{itemize}

%--------------------------------------------------------------------
\section{Domande d'orale probabili, con risposta breve}

\begin{enumerate}
  \item \textbf{Cos'è il double descent?} Il test error, in funzione della complessità, scende, risale fino a un picco
        alla soglia di interpolazione ($d=m$) e poi scende di nuovo nel regime sovra-parametrizzato.
  \item \textbf{Perché il picco proprio a $d = m$?} Lì $X$ è quasi quadrata, $X^\top X$ ha un autovalore quasi nullo,
        la soluzione divide il rumore per quell'autovalore e $\norm{w}$ esplode (slide 9, p.~5 e Sezione 5).
  \item \textbf{Perché il training error è $0$ per $d \ge m$?} Con $d \ge m$ e righe indipendenti, $Xw = y$ ha almeno
        una soluzione esatta: il modello interpola.
  \item \textbf{Che soluzione usi per $d > m$ e perché?} La minima norma $X^{+}y$: è unica, interpola, ed è il limite
        della ridge per $\alpha \to 0$ e il punto a cui converge la GD partita da $0$.
  \item \textbf{Come hai calcolato $X^{+}$?} Con \texttt{np.linalg.pinv}, basata sulla SVD. È algebra lineare, non un fit.
  \item \textbf{Cosa fa la ridge al picco?} Sostituisce $\lambda_i$ con $\lambda_i + \alpha$ al denominatore:
        niente amplificazione del rumore, picco attenuato o eliminato, al costo di un po' di bias.
  \item \textbf{Come hai scelto $\alpha$?} Valori fissi a confronto, e/o validazione solo sul training set. Mai sul test.
  \item \textbf{Come fai variare $d$ sul dataset reale?} Con feature casuali; il modello resta lineare nei parametri.
  \item \textbf{Come eviti il data leakage?} Split prima di tutto, statistiche di standardizzazione solo dal training,
        iperparametri solo dal training.
  \item \textbf{Perché il double descent contraddice la teoria classica?} La teoria classica prevede che oltre un certo
        punto aumentare la complessità peggiori sempre la generalizzazione (curva a U). Nel regime $d > m$ succede
        il contrario.
  \item \textbf{Punti deboli del tuo lavoro?} Modelli solo lineari; il dataset reale richiede feature casuali, quindi il
        risultato dipende da come le costruisci; poche ripetizioni per vincoli di tempo; scelta di $m$ arbitraria.
\end{enumerate}

%--------------------------------------------------------------------
\section{Soluzioni delle domande di verifica}
\emph{Prova a rispondere prima da solo, poi confronta.}

\paragraph{Domanda 1.} Perché per $d > m$ la formula $(X^\top X)^{-1}X^\top y$ non si può usare, e cosa calcola \texttt{pinv(X) @ y}?

\smallskip
\noindent\textit{Risposta.} $X^\top X$ è $d\times d$ ma ha rango al massimo $m < d$, quindi è singolare e non invertibile
(equivalentemente, $m$ vettori non generano $\R^d$, slide 9 p.~4). Il sistema $Xw = y$ ha infinite soluzioni che
interpolano. \texttt{pinv(X) @ y} restituisce quella di norma euclidea minima, $X^\top(XX^\top)^{-1}y$: è unica, sta
nello spazio delle righe di $X$, e ogni altra soluzione ha norma maggiore (Pitagora). Per $d \le m$ coincide con la
formula delle slide.

\paragraph{Domanda 2.} Perché la ridge non ha il picco a $d = m$?

\smallskip
\noindent\textit{Risposta.} Nella soluzione di minima norma le componenti di $y$ vengono divise per gli autovalori
$\lambda_i$ di $X^\top X$; a $d \approx m$ il più piccolo è quasi $0$ e il rumore esplode. La ridge inverte
$\alpha I + X^\top X$, i cui autovalori sono $\lambda_i + \alpha \ge \alpha > 0$ (slide 9 p.~5): il denominatore non
scende mai sotto $\alpha$, quindi $\norm{w}$ e il test error restano limitati. Il prezzo è il bias: con $\alpha$ troppo
grande $w \to 0$ e il modello sotto-adatta.

%--------------------------------------------------------------------
\section{Fonti}
\begin{itemize}
  \item N.\ Cesa-Bianchi, slide ``Introduction'' (1.\ intro), p.~3--5; ``Statistical learning'' (4.\ statRisk), p.~1--6;
        ``Hyperparameter tuning and risk estimates'' (6.\ crossVal), p.~2--4; ``Kernels'' (11.\ kernels), p.~1.
  \item N.\ Cesa-Bianchi, slide ``Linear predictors'' (9.\ linear), versione del 10/02/2026, p.~4--5.
  \item N.\ Cesa-Bianchi, slide ``Online Gradient Descent'' (10.\ online), versione del 24/02/2026.
  \item M.\ Belkin, D.\ Hsu, S.\ Ma, S.\ Mandal, ``Reconciling modern machine learning practice and the classical
        bias-variance trade-off'', PNAS 2019 (arXiv:1812.11118). Versione arXiv v2 (10/09/2019); numeri di pagina riferiti a quel PDF.
  \item M.\ Belkin, D.\ Hsu, J.\ Xu, ``Two models of double descent for weak features'', 2019 (arXiv:1903.07571). \emph{Esterna.}
  \item T.\ Hastie, A.\ Montanari, S.\ Rosset, R.\ Tibshirani, ``Surprises in high-dimensional ridgeless least squares
        interpolation'', Annals of Statistics 2022. \emph{Esterna.}
  \item P.\ Nakkiran, P.\ Venkat, S.\ Kakade, T.\ Ma, ``Optimal regularization can mitigate double descent'', ICLR 2021. \emph{Esterna.}
  \item Assignment ``ML projects 25/26'', Assignment 4 (p.~8--9) e FAQ (p.~16).
\end{itemize}

\end{document}
